% Fichier généré automatiquement par tools/generate_corrections.py.
% Source : DYN/DYN-04-TorseurDynamique/06_TR/06_TR.tex
% Statut : complété (3/3 question(s)).
\begin{corrigebox}[Corrigé complété]
\CorrigeQuestion{1}
\textbf{Expression de la résultante dynamique}
$\vectrd{2}{0}=m_2\vectg{G_2}{2}{0} = m_2\dderiv{\vect{AC}}{\rep{0}}$
$\dderiv{\vect{AC}}{\rep{0}}= \dderiv{\vect{AB}}{\rep{0}}+\dderiv{\vect{BC}}{\rep{0}}$
$= \ddot{\lambda}(t) \vect{i_0} + R \dderiv{\vi{2}}{\rep{0}}$
$= \ddot{\lambda}(t) \vect{i_0} + R \deriv{\dot{\theta} \vj{2}}{\rep0}$
$= \ddot{\lambda}(t) \vect{i_0} + R \left(\ddot{\theta} \vj{2} - \dot{\theta}^2 \vi{2}\right) $.

\textbf{Méthode 1 : Calcul en $G_2=C$ puis déplacement du torseur dynamique}
\begin{itemize}
\item Calcul du moment cinétique en $G_2$ : $G_2=C$ est le centre de gravité donc $\vectmc{C}{2}{0}=\inertie{C}{2}\dot{\theta}\vect{k_0} = C_1  \dot{\theta} \vk{1}$.
\item Calcul du moment dynamique en $G_2$ : $G_2=C$ est le centre de gravité donc $\vectmd{C}{2}{0}=\deriv{\vectmc{C}{2}{0}}{\rep{0}}= C_1  \ddot{\theta} \vk{1}$.
\item Calcul du moment dynamique en $B$ : $\babard{B}{C}{2}{0} $
$= C_2  \ddot{\theta} \vk{1} + R\vi{2}m_2\wedge\left(  \ddot{\lambda}(t) \vect{i_0} + R \left(\ddot{\theta} \vj{2} - \dot{\theta}^2 \vi{2}\right)\right)$
$= C_2  \ddot{\theta} \vk{1} + Rm_2\left( -\sin \theta \ddot{\lambda}(t) \vk{0} + R \ddot{\theta} \vk{2}\right)$.
\end{itemize}

Au final, on a donc $\torseurdyn{2}{0}=\torseurl{
m_2\left(\ddot{\lambda}(t) \vect{i_0} + R \left(\ddot{\theta} \vj{2} - \dot{\theta}^2 \vi{2}\right)\right)
}{
C_2  \ddot{\theta} \vk{1} + Rm_2\left( -\sin \theta \ddot{\lambda}(t) \vk{0} + R \ddot{\theta} \vk{2}\right)
}{B}$.
\CorrigeQuestion{2}
On a $\vectrd{1+2}{0}=\vectrd{1}{0}+\vectrd{2}{0}=m_1\ddot{\lambda}(t)\vi{0}+m_2\left(\ddot{\lambda}(t) \vect{i_0} + R \left(\ddot{\theta} \vj{2} - \dot{\theta}^2 \vi{2}\right)\right)$.
On projette alors sur $\vect{i_0}$,
$\vectrd{1+2}{0}\cdot \vect{i_0} =m_1\ddot{\lambda}(t)+m_2\left(\ddot{\lambda}(t)-R \left(\ddot{\theta} \sin\theta(t)  + \dot{\theta}^2 \cos\theta \right)\right) $.
\CorrigeQuestion{3}
On choisit les inconnues et les conventions positives de la figure,
        on écrit les relations de conservation/fermeture adaptées, puis on résout le
        système obtenu. Le résultat symbolique doit être homogène et vérifié dans les cas
        limites (entrée nulle, rapport unitaire ou position de référence).
\end{corrigebox}
